\section{Comparative Analysis}\label{sec:comp}

\subsection{Overall Performance Comparison}

To show how the usual evaluation differs from realistic conditions, the $R^2$ of five
model--strategy combinations is compared in Table~\ref{tab:prior} under the random split with
complete data, which is the protocol of most earlier studies, and under four realistic
conditions: the temporal split, the district holdout, and the district holdout with 30\% of the
weather missing under each pattern. The last column gives the gap between the random-split $R^2$
and the mean $R^2$ under the three missingness patterns in unseen districts.

\begin{table}[!t]
\caption{$R^2$ Under the Usual Random-Split Protocol and Under Realistic Conditions (Mean Over
Crops; Best per Column in Bold)}\label{tab:prior}
\centering\footnotesize\setlength{\tabcolsep}{2.4pt}
\begin{tabular}{lrrrrrrr}
\toprule
 & Usual & \multicolumn{2}{c}{Complete data} & \multicolumn{3}{c}{Distr.\ + 30\% missing} & \\ \cmidrule(lr){3-4}\cmidrule(lr){5-7}Strategy & random & Temp. & Distr. & MCAR & Var. & Block & Gap \\
\midrule
Ridge + median & 0.31 & \textbf{0.04} & 0.19 & 0.13 & 0.14 & 0.17 & 0.17 \\
XGB + median & 0.32 & $-$0.04 & 0.24 & 0.20 & 0.17 & 0.20 & 0.13 \\
XGB + native & 0.32 & $-$0.04 & 0.24 & 0.19 & 0.18 & 0.21 & 0.13 \\
RF + median & \textbf{0.37} & $-$0.03 & \textbf{0.29} & 0.25 & \textbf{0.23} & \textbf{0.25} & 0.13 \\
RF + iterative & \textbf{0.37} & $-$0.03 & \textbf{0.29} & \textbf{0.26} & 0.22 & 0.25 & 0.12 \\
\bottomrule
\end{tabular}


\vspace{2pt}{\scriptsize Temp.: temporal split; Distr.: district holdout; Var.: variable-wise;
Gap: random-split $R^2$ minus the mean $R^2$ of the three missingness columns.}
\end{table}

The random split overstated the accuracy available in unseen districts with incomplete weather by
0.12 to 0.17 in $R^2$ for every strategy, that is, by about a third to a half of the reported
value. The random forest with iterative imputation, which is recommended by the results of
Section~\ref{sec:results}, retained the highest accuracy under random masking ($R^2 = 0.26$) and
shared the best values under the other patterns with the random forest with median replacement.
Compared with the strategies most often found in earlier work, it retained 0.07 more $R^2$ than
XGBoost with native handling under random masking and 0.13 more than ridge regression with median
replacement, while the gap to its own random-split value was the smallest of all strategies (0.12).
These gains were obtained with existing methods; the comparison is therefore presented as evidence
for a choice among them, not as a new model, and because the strategies were compared on the
same data on which the recommendation is based, the advantage should be confirmed in other states.

\subsection{Comparison Among the Models}

A consistent picture is formed when the results are taken together. The random forest was the
most accurate model under the random split and the district holdout, with complete and with
incomplete weather; a random-forest strategy was the best of the ten in 35 of the 36 combinations
of crop, pattern and rate, regardless of the imputer. XGBoost was second, and its
native handling of missing values was not better than imputation. Ridge regression was the least
damaged by missingness, the best model under the temporal split and, with location features, the
best model on unseen districts with complete data; its strength lies in extrapolation, which tree
ensembles cannot provide. The mean predictor was clearly outperformed in every design, but all
models explained at most about 40\% of the yield variance, which reflects the limited information
in district-level weather summaries.

\subsection{Comparison With Earlier Studies}

The main findings are set against earlier work in Table~\ref{tab:earlier}. The over-optimism of
random validation reported for spatial data \cite{roberts2017,ploton2020} and for yield models in
the United States and Africa \cite{chakravarty2026,adjei2026} was confirmed for Indian districts,
although the loss was smaller: positive $R^2$ was retained in unseen districts of one state, where
neighbouring districts share climate and practices. As in Germany \cite{baatz2025}, transfer to
future years was harder than transfer to new locations. The observation that no single imputer is
always best \cite{saini2024,calin2023} was confirmed, and it was found to matter less than the
choice of model.

\begin{table}[!t]
\caption{Findings of This Study Compared With Earlier Work}\label{tab:earlier}
\centering\footnotesize\setlength{\tabcolsep}{2.5pt}
\begin{tabular}{p{1.35cm}p{2.05cm}p{2.1cm}p{1.8cm}}
\toprule
Study & Earlier finding & Result in this study & Added insight \\
\midrule
Alsaber \emph{et al.} \cite{alsaber2025} & No single model best; overfitting on test years & An RF strategy was best in 35 of 36 settings of unseen districts & Ranking is stable once evaluation is realistic \\
Manogna R.~L. \emph{et al.} \cite{manogna2026} & Persistence nearly ties a deep model & Weather explains at most 40\% of variance & Gains over naive baselines are modest \\
Saini \& Nagpal \cite{saini2024} & Mean replacement competitive & Median within 0.4 pp of iterative & Model choice matters more \\
Chakravarty \cite{chakravarty2026}, Adjei \cite{adjei2026} & Negative $R^2$ in unseen regions & $R^2$ falls from 0.37 to 0.29 & Smaller loss within one state \\
Baatz \cite{baatz2025} & Temporal transfer harder than spatial & Confirmed; linear trend helps & Trees cannot extrapolate trends \\
\bottomrule
\end{tabular}
\end{table}

\subsection{Practical Implications}

A simple guide, summarised in Table~\ref{tab:guide}, is suggested by the findings. A model that
will be used in districts outside its training data should be evaluated with a district holdout,
because the random split overstates its accuracy. A random forest is the safest model under
geographic shift and missing weather, and iterative imputation should be preferred when values
are missing at random, although median replacement is nearly as good. When whole variable groups
are missing, no imputer recovers the loss, and the variable should be obtained from another
source. When the model must predict future seasons, a model with a linear trend term, or a
detrended target, is needed. Because missingness and geographic shift did not compound, a model
that has been validated on unseen districts can be expected to tolerate moderate gaps in its
weather inputs.

\begin{table}[!t]
\caption{Practical Guide Derived From the Results}\label{tab:guide}
\centering\footnotesize\setlength{\tabcolsep}{3pt}
\begin{tabular}{p{2.2cm}p{2.6cm}p{2.6cm}}
\toprule
Situation & Main finding & Recommendation \\
\midrule
New districts & $R^2$ drops by 0.08--0.13 & Validate with district holdout; use RF \\
Random gaps & Small damage (0.4--0.8 pp at 30\%) & RF with iterative or median imputation \\
Missing variable group & Largest damage; imputation does not help & Obtain the variable from another source \\
Outages in daily data & Features distorted, not missing & Treat as noise; any imputer \\
Future seasons & Trees fail to extrapolate & Ridge with trend, or detrended target \\
\bottomrule
\end{tabular}
\end{table}
